% Section of golfer_appearance_reference.tex (R8, #12685). Standalone research
% and visual reference, not the design manual.
\section{One Shared Appearance and Anthropometry Spec (R8, \#12685)}

\subsection{Defect}
Skin, garment, blend-ball and shoe meshes were built only inside the MuJoCo
appearance layer (\texttt{appearance\_layer.\_body\_meshes}), so every other
engine drew grey capsules. The trunk had three independent size sources: the
spec \texttt{visual\_hints} ellipsoids (stature fractions
\texttt{TORSO\_WIDTH\_FRACTION}/\texttt{TORSO\_DEPTH\_FRACTION}), the fixed
\texttt{library.PART\_RADIUS\_M} metres with \texttt{PART\_ASPECT}, and the
GS3DX absolute radii. The MuJoCo trunk therefore did not scale with the
subject, and the trunk body carried three nearly identical overlapping lofts
(to the cervicale and both shoulder hubs). The grey ``plain'' palette was
copied into four modules. The default outfit was navy polo over khaki shorts
whose thigh band covered only 60\,\% of the thigh, leaving a bare-knee band.

\subsection{Contract}
\texttt{model\_appearance.body\_meshes} is pure data plus pure functions:
\begin{itemize}
\item \texttt{body\_meshes(spec, doc, club=None)} $\to$
  \texttt{list[BodyMesh(body, label, mesh\_in\_body\_frame, material\_name)]}:
  skin (\texttt{skin}$i$), garment (\texttt{garment}$i$), blend ball
  (\texttt{blend}$i$) and shoe meshes, plus the club parts when an assembly is
  given. Deterministic; bodies in visual-skeleton order; the head body is
  omitted while the parametric head (\texttt{model\_appearance.head}) is on.
\item \texttt{resolve\_appearance(spec, doc, club=None, head\_mesh=None)} $\to$
  \texttt{ResolvedAppearance(meshes, materials, head\_mesh, frame)}: the only
  object an engine emitter receives (Law of Demeter). \texttt{materials} maps
  every referenced name to its PBR \texttt{Material};
  \texttt{rgba(name)} and \texttt{meshes\_for(body)} delegate.
\item \texttt{section\_size(spec, body)} $\to$
  \texttt{SectionSize(radius\_m, aspect, source)} for a trunk body.
\item \texttt{MeshReference(path, licence, provenance)}: the reserved head
  slot for R7 (\#12687); every field must be non-empty.
\end{itemize}
Preconditions raise \texttt{ValueError}: an unknown material name (document
segment rule, clothing preset or a produced mesh), a hinted trunk body without
landmark axes, a trunk aspect outside $(0,1]$ (a trunk deeper than wide is a
spec error). A body that matches no part rule classifies as \texttt{other} and
gets the plain grey material \texttt{plain\_grey}, never a skin tone.

\subsection{Frames}
Every \texttt{mesh\_in\_body\_frame} vertex is in the \emph{spec body frame}
(\texttt{SPEC\_BODY\_FRAME = "spec\_body"}): the frame the spec joints,
\texttt{visual\_skeleton} capsules and \texttt{visual\_hints} shapes use. An
engine whose body frame differs transforms explicitly. The MuJoCo exporter's
body is the joint follower frame, so the emitter maps a spec point $p$ to
$T_b\,p$ with $T_b = (\texttt{child\_to\_follower})^{-1}$ (the exporter's
\texttt{offsets[body]}). The v1 chest body \texttt{COMRod} has
$R_{\mathrm{c2f}} = \left(\begin{smallmatrix}0&0&1\\1&0&0\\0&1&0\end{smallmatrix}\right)$;
attaching spec-frame geometry without $T_b$ misorients the chest. The
regression test compiles the MJCF in MuJoCo and checks the chest mesh in the
\emph{world} frame (\texttt{geom\_xpos} $+$ \texttt{geom\_xmat}$\cdot$
\texttt{mesh\_vert}) against the world shoulder line; replacing $T_b$ by a
pure translation fails it for v1.

\subsection{Pelvis Solid (HRG-C1, \#13164)}
A pelvis body that parents two hips and a torso joint is not lofted on a
capsule: after a functional hip calibration the longest capsule was the left
hip and the pelvis was drawn $51$\,mm off the midline. It is instead a convex
split ellipsoid. With up axis $\hat u$ (\texttt{BodySection.trunk\_up\_axis}),
hip joint centres $p_r,p_l$ and lumbar joint $p_L$ (all in the pelvis frame),
the centre is $c = \bar p - (\bar p\!\cdot\!\hat u)\hat u + z_{\min}\hat u$ with
$\bar p=(p_r+p_l)/2$, $z_{\min}=\min_i p_i\!\cdot\!\hat u$ (widest level). The
solid runs $\delta=r_{\mathrm{thigh}}$ below $c$ and is cut flat at the lumbar
plane, height $T = p_L\!\cdot\!\hat u - z_{\min}$; the upper half-height is
$T/\sqrt{1-\kappa^2}$ with waist ratio $\kappa=0.75$. Half width
$W=\max\bigl(\rho,\ \max_i (|(p_i-\bar p)\!\cdot\!\hat m| + r_{\mathrm{thigh}} + 0.01)
/\sqrt{1-\ell_i^2}\bigr)$ with $\ell_i$ the HJC's height over the widest level
in upper half-heights; half depth is the hint AP semi-axis. The chest's lower
tip is its joint origin, so the two solids meet at the lumbar plane without
overlap. The thigh blend ball is $0.95\,r_{\mathrm{thigh}}$ ($\le1.05$ with the
garment). Visual only: compiled masses, inertias, joints and equality
constraints are unchanged (\texttt{test\_pelvis\_loft.py}, the compiled-array
test in \texttt{test\_mujoco\_appearance\_layer\_body\_meshes.py}).

\subsection{Pelvis Width (HRG-C4, \#13187)}
C1's width rule enclosed each whole thigh radius, $W=\max_i(|(p_i-\bar p)\!\cdot\!\hat m|
+ r_{\mathrm{thigh}}+0.01)/\sqrt{1-\ell_i^2}$, which gave a $395$\,mm skin and
$427$\,mm trouser width on the fitted best-driver capture (and $354$\,mm at the
corrected $178$\,mm hip spacing). Only the femoral heads belong inside the
pelvis; the thigh joins through the hip ball. The \emph{full} skin width
$2W$ (at the widest level, which is where the section correction is $1$) is now
\begin{equation}
2W=\max\bigl(B,\ 2\max_i\frac{|(p_i-\bar p)\cdot\hat m|+r_{\mathrm{head}}+m}
{\sqrt{1-\ell_i^2}}\bigr),\quad r_{\mathrm{head}}=25\,\mathrm{mm},\ m=10\,\mathrm{mm},
\end{equation}
with target $B$ chosen in this order: the subject's measured standing hip
breadth (\texttt{subject.hip\_breadth\_m}, source \texttt{measured}); else
$B=f_H\,H$ with stature $H$ (\texttt{subject.stature\_m}, source
\texttt{stature\_prior}); else the spec hint ML width (\texttt{spec\_hint}).
When the containment term wins, the width follows it and the status is
\texttt{containment\_dominated}; it is never clipped
(\texttt{body\_meshes.pelvis\_width\_report}). Because the HJC term is
$0.089+0.035=0.124$\,m per side at $178$\,mm and $0.110+0.035=0.145$\,m at $220$\,mm,
both stay below $B/2$, so the result is identical for the two spacings and a
later HJC correction (C3, \#13172) does not retune the visual. The lower pole is
$\delta=r_{\mathrm{thigh}}$ below the \emph{mean} HJC height $\bar z$ (so vertical
HJC asymmetry does not stretch the solid), with at least $25$\,mm of solid below
the lowest HJC; the lumbar interface and the C1 gates (symmetry $<1$\,mm,
centring $<5$\,mm, lumbar overlap $\le10^{-6}$\,m) are unchanged.

\begin{table}[h]
\centering
\begin{tabular}{llll}
Quantity & Value & Status & Source \\ \hline
$f_H$, hip breadth / stature & $0.195$ ($0.333$\,m at $1.71$\,m) & \textbf{unverified prior} & panel working figure \\
ANSUR II male standing hip breadth & mean and SD not obtained & unverified & Gordon et al.\ 2014 (not accessible here) \\
Femoral head radius & $25$\,mm & modelling choice & --- \\
Containment margin & $10$\,mm & modelling choice & --- \\
\end{tabular}
\end{table}
The $0.195$ prior could not be checked against the published ANSUR II table
(the report PDF returned HTTP 403 and no search result gave the mean), so it
is recorded as an unverified prior, not a measurement; the ANSUR II public CSV
(\texttt{hipbreadth}, male) should replace it with a mean and SD. Supplying
\texttt{subject.hip\_breadth\_m} overrides it. Measured result at $H=1.71$\,m:
skin $333.5$\,mm (target $333.45$\,mm), trousers $360.1$\,mm, versus
$395.2$ / $426.8$\,mm before. Visual only; physics unchanged. Tests:
\texttt{tests/unit/model\_appearance/test\_pelvis\_width.py}. Reproduce stills:
render the best-driver bundle before and after at address, top and impact.

\subsection{Sizing Rule}
Limbs, hands, feet and the head keep \texttt{PART\_RADIUS\_M} (with the
short-link clamp $\rho\le 0.6\,\ell$). The trunk (torso, pelvis) is one lofted
solid per body along its longest capsule (joint origin to the farthest child
joint). With the body's landmark axes $\hat m$ (ML) and $\hat d$ (AP) from
\texttt{cross\_section.body\_sections} and its hint ellipsoids
$\{h_k\}$ (axis-aligned half sizes in the spec body frame), the semi-axis of
an ellipsoid along a unit $\hat e$ is $r_k(\hat e) = \bigl(\sum_i (e_i/h_{k,i})^2\bigr)^{-1/2}$
(a box uses $\sum_i |e_i| h_{k,i}$), and
\[
\rho = \max_k r_k(\hat m), \qquad \alpha = \frac{\max_k r_k(\hat d)}{\rho},
\qquad \texttt{source} = \texttt{visual\_hints}.
\]
The loft extents are then $2\rho$ along $\hat m$ and $2\alpha\rho$ along
$\hat d$ (the R1 ring formula). Hint sizes are anthropometric, so no
length clamp applies. A trunk body without hint shapes (the v1 spec; the
anthro \texttt{UpperTorsoBase} stub) uses $\rho=\texttt{PART\_RADIUS\_M}$,
$\alpha=\texttt{PART\_ASPECT}$, \texttt{source = library\_default}, with the
clamp. For the anthropometric captures (stature $1.71$\,m) the thorax hint
gives $\rho = 0.185\times1.71/2 = 0.1582$\,m (a \emph{half} width; full
shoulder-line chest width $0.316$\,m) and $\alpha = 0.6486$.

\subsection{Outfit and Coverage}
Owner ruling (\#12685): the default outfit is charcoal trousers.
\texttt{schema.DEFAULT\_CLOTHING = "golf\_polo\_trousers"} (white polo, black
shoes) and \texttt{library.TROUSERS\_MATERIAL = "trousers\_charcoal"} is the one
trouser token the pelvis and both thighs and shins wear. Coverage is per
outfit: \texttt{garment\_coverage(preset, part)} returns $(0,1)$ for every leg
segment of a full-length preset (its shins wear its pants), so each thigh and
shin garment is a closed spindle that encloses the skin spindle; shorts keep
the $(0,0.6)$ thigh band. The test checks, at the zero pose in world
coordinates, that every leg skin vertex lies inside its garment hull and that
the vertical garment spans pelvis $\to$ thigh $\to$ shin $\to$ shoe overlap.

\subsection{One Plain Palette}
\texttt{library.PLAIN\_PALETTE} holds the capsule, hint-shape, COM, frame,
floor and ball RGBA; \texttt{plain\_rgba(key)} validates the key. The MuJoCo
\texttt{visual\_layer} reads it; Drake, Pinocchio and OpenSim adopt it in
R9a/R9b.

\subsection{Engine Hooks}
MuJoCo (\texttt{appearance\_layer.attach\_appearance}) is the reference
consumer: it calls \texttt{resolve\_appearance} and only writes MJCF assets and
geoms. \texttt{native\_viewer\_export.backends.\_body.body\_mesh\_files(resolved,
directory)} writes one OBJ per \texttt{BodyMesh} in its spec body frame with
its library RGBA (\texttt{BodyMeshFile(body, label, path, material\_name,
rgba)}), mirroring \texttt{\_head.py}; the MeshCat and OpenSim backends register
these files in R9a/R9b (\#12688, \#12689), and Simscape reads the same
resolution in R9c (\#12690).

\subsection{Before and After}
Spec-frame skin extents (ML $\times$ AP, metres) for both anthropometric captures
(identical trunk geometry):
\begin{center}
\begin{tabular}{lll}
Body & Before (main after R1/R2) & After (R8) \\ \hline
\texttt{COMRod} (chest) & 3 lofts, $0.290\times0.209$ & 1 solid, $0.316\times0.205$ \\
\texttt{LowerTorso} (pelvis) & 3 lofts, $0.166\times0.129$ (clamped) & 1 solid, $0.290\times0.196$ \\
Hint target $2\rho\times2\alpha\rho$ & & $0.316\times0.205$; $0.291\times0.197$ \\
\end{tabular}
\end{center}
Every non-trunk body is unchanged: a golden per-body summary (material, vertex
count, MJCF-frame bounding box) captured from the pre-refactor layer with the
shorts preset matches to $2\times10^{-6}$\,m. Physics is bit-identical (masses,
inertias, body positions, joint axes). Address stills face-on and
down-the-line, capture-A driver and capture-B 7-iron, are in
\texttt{evidence/appearance/shared\_appearance\_spec/} (after) beside the R1
stills in \texttt{evidence/appearance/trunk\_cross\_section/} (before, navy
polo over khaki shorts with a bare knee): the after stills show charcoal
trousers hip to ankle, a white polo, no blue segment and a chest wider than
deep. \texttt{trunk\_and\_outfit.json} records the sizing source and garment
of each trunk and leg body. Stills are visual evidence and do not qualify
dynamics.

\subsection{Limitations}
\begin{itemize}
\item The skinned glTF body (\texttt{skinned\_body}) still sizes its trunk from
  \texttt{PART\_RADIUS\_M}; migrating it to \texttt{section\_size} is follow-up
  work outside this change's file set.
\item The anthro \texttt{UpperTorsoBase} is a $6$\,mm stub without hints; its
  library-default solid is clamped to $\rho = 0.6\,\ell$ and is invisible.
\item GS3DX absolute radii are unchanged until R9c reads this resolution.
\item The trunk solid follows the longest capsule, so a spec whose spine is
  not its longest trunk capsule would need an explicit spine landmark.
\end{itemize}

\subsection{Reproduction}
\begin{verbatim}
python -m pytest -o addopts="" tests/unit/model_appearance/test_body_meshes.py \
  tests/unit/motion_matching/test_mujoco_appearance_layer_body_meshes.py \
  tests/unit/tools/native_viewer_export/test_body_mesh_files.py -q
PYTHONPATH=. python3 docs/development/full_body_models/evidence/appearance/\
shared_appearance_spec/render_shared_appearance_spec.py \
  --datastore D:/FleetData/swing-datastore --out OUT --evidence \
  docs/development/full_body_models/evidence/appearance/shared_appearance_spec
\end{verbatim}
